\documentclass[10pt,a4paper]{amsart}
\usepackage[margin=19mm]{geometry}
\usepackage{amsmath,amssymb,mathtools}
\usepackage[scr=rsfs,cal=euler]{mathalfa}
\usepackage{xfrac}
\usepackage[unicode,hidelinks,bookmarks=false]{hyperref}
\newcommand{\ie}{i.e.\ }
\newcommand{\cf}{cf.\ }
\newcommand{\resp}{respectively\ }
\newcommand{\Subsection}{\subsection}
\providecommand{\nobreakdash}{\nobreak\hbox{-}}
\setlength{\emergencystretch}{2em}
\multlinegap0pt
\usepackage{amssymb}
\usepackage[shortlabels]{enumitem}
\newtheorem{lemma}{Lemma}
\newcommand{\N}{\mathbb{N}}
\newcommand{\Z}{\mathbb{Z}}
\title{On the chain of commuting operators on Banach spaces}
\author{Tomasz Szczepanski}
\date{Source: arXiv:2507.14297v2 (CC BY 4.0)}
\begin{document}
\maketitle

\noindent\emph{Source and licence.} On the chain of commuting operators on Banach spaces. Version arXiv:2507.14297v2 (CC BY 4.0), distributed by arXiv under the Creative Commons Attribution 4.0 International licence, http://creativecommons.org/licenses/by/4.0/, as recorded in the arXiv licence metadata for this exact version and shown on the abstract page https://arxiv.org/abs/2507.14297v2. Full licence notice: \texttt{licenses/arxiv-2507.14297-CC-BY-4.0.txt}.\par\smallskip

\noindent\textit{Editorial scope.} Counted unit: the article's lemma sigmaset (source0.tex lines 289--293). The article's complete original proof of it is delivered with it (source0.tex lines 294--321, one \texttt{proof} environment of 28 lines). No mathematics is written, repaired or re-derived: the statement and the proof are the article's own lines, in the article's own notation, and no retained line is substituted or re-flowed.  No further source-local context is needed: every cross-reference of every retained line resolves inside the two delivered spans.  Nothing was omitted from any retained span, so the package carries no omission record.  No retained line contains a citation, so the wrapper carries no bibliography and no bibliography entry.  No retained line contains a figure environment, an image-inclusion command or drawing code, so no image is delivered and the package declares no asset dependency.  Editorial formatting only, in addition: an AMS \texttt{amsart} preamble that restates the article's own declarations and macros used by the retained lines, namely the environments \texttt{lemma} on separate counters, together with the standard packages those lines need.  The retained environments print in this document as Lemma 1 in order of appearance, whereas the article prints the same items with the numbering of its own full document.  The complete native source of the article as distributed in the arXiv e-print for this exact version is bundled unchanged as \texttt{sources/181832/source0.tex}.
\begin{lemma}
\label{sigmaset}
Let $\sigma:\N \rightarrow \N$ be injective. There exists a proper, non-empty subset $B \subsetneq \N$ with the following property: for every $n\in \N$,
$$n \in B \iff \sigma^2(n)\in B.$$ 
\end{lemma}

\begin{proof}
We consider two separate cases.

Case 1. $\sigma$ is onto. Fix $n_0 \in \N$ and consider the sets
$$A = \{\sigma^k(n_0): k\in \Z\} ~\text{ and } ~B = \{\sigma^{2k}(n_0): k \in \Z\}.$$
First, note that $\varnothing \neq B \subseteq A \subseteq \N$. Our goal is to show that $B$ is a proper subset of $\N$. If $A$ is a proper subset of $\N$, then automatically so is $B$. Let us assume that $A=\N$. Suppose, by way of contradiction, that $B=\N=A$. This implies that $\sigma(n_0) \in B$, so there exists a $k\in \Z$ such that $\sigma(n_0) = \sigma^{2k}(n_0)$. But then
\begin{align*}
A &= \{\sigma^{-2k+1}(n_0),\sigma^{-2k+2}(n_0),\dots,\sigma^{-1}(n_0),n_0, \sigma(n_0),\dots,\sigma^{2k-1}(n_0)\}. 
\end{align*}
In particular $A$ is finite, yielding a contradiction. Hence $B\subsetneq\N$. Moreover, it is straightforward to check that 
$$n \in B \iff \sigma^2(n)\in B,$$
which concludes Case 1.

Case 2. $\sigma$ is not onto. Just like in Case 1, we fix $n_0 \in \N$ but we additionally assume that $n_0\notin \operatorname{Range} (\sigma)$. Consider the sets 
$$A = \{\sigma^k(n_0): k\in \N\cup\{0\}\} ~\text{ and } ~B = \{\sigma^{2k}(n_0): k \in \N\cup \{0\}\}.$$
We will show that $B$ satisfies the claim. To show $B\neq \N$, following the same argument as in Case 1, suppose, by way of contradiction, that $B=\N=A$. Once again, since $\sigma(n_0) \in B$, we can find $k\in \N\cup \{0\}$ such that $\sigma(n_0) = \sigma^{2k}(n_0)$. But then
\begin{align*}
A =\{n_0,\sigma(n_0),\dots,\sigma^{2k-1}(n_0)\},
\end{align*}
which again results in a contradiction. Hence $B\subsetneq\N$. 

Finally we show that
$$n \in B \iff \sigma^2(n)\in B.$$ 
First, suppose $n\in B$. We can find $k\in \N\cup\{0\}$ such that
$n=\sigma^{2k}(n_0)$. Thus, $\sigma^2(n)=\sigma^{2(k+1)}(n_0) \in B$. 

Suppose now $n$ is such that $\sigma^2(n) \in B$. There exists a $k \in \N\cup\{0\}$ with $\sigma^2(n)=\sigma^{2k}(n_0)$. Note that $k\neq 0$, as this would imply $n_0=\sigma^2(n)\in \operatorname{Range}(\sigma^2) \subseteq \operatorname{Range}( \sigma)$, which contradicts $n_0 \notin \operatorname{Range}(\sigma)$. Therefore, $k\geq 1$ and $n=\sigma^{2k-2}(n_0) \in B$ follows.
\end{proof}

\end{document}
